%------------------------------------------------------------------------------%
\begin{exercises}

% Original
%------------------------------------------------------------------------------%
\begin{exercise}
  Escreva as quatro primeiras somas parciais $S_n$ da série.
  \begin{mathlist}[2]
    \item \sum_{n=1}^\infty (-1)^n\;n
    \item \sum_{n=0}^\infty \frac{n}{2}
    \item \sum_{n=0}^\infty 2^{2-n}
    \item \sum_{n=1}^\infty \ln(n)
  \end{mathlist}
\end{exercise}

% Original
%------------------------------------------------------------------------------%
\begin{exercise}
  Escreva a série como um somatório.
  \begin{mathlist}
    \item 1 + 3 + 5 + 7 + 9 + 11 + \cdots
    \item 2 - 4 + 6 - 8 + 10 - 12 + \cdots
    \item 3
      + \frac{5}{2}
      + \frac{7}{4}
      + \frac{9}{8}
      + \frac{11}{16}
      + \cdots
    \item 1
      + 0,1
      + 0,01
      + 0,001
      + 0,0001
      + \cdots
  \end{mathlist}
\end{exercise}

% Seção 10.2 - Exercício 2 página 20 Thomas
%------------------------------------------------------------------------------%
\begin{exercise}
Encontre a fórmula para a $n$-ésima soma parcial da série
\[
  \frac{9}{100} +
  \frac{9}{100^2} +
  \frac{9}{100^3} + \cdots
  \frac{9}{100^n} + \cdots
\]
use-a para encontrar a soma da série se ela convergir.
\begin{answer}
  Observamos que os termos da série são
  \[
    a_n = \frac{9}{100^n} = \frac{9}{100} \left(\frac{1}{100}\right)^{n-1}
  \]
  ou seja, essa é uma série geométrica com
  \[
    a = \frac{9}{100} \qquad \text{ e } \qquad r = \frac{9}{100}
  \]
  a fórmula para as somas parciais de uma série geométrica é
  \[
    S_n = \frac{a(1-r^n)}{1-r}
  \]
  substituindo os valores de $a$ e $r$ temos
  \[
    S_n = \dfrac{\dfrac{9}{100}\left[1-\left(\dfrac{1}{100}\right)^n\right]}{1-\dfrac{1}{100}}
  \]
  calculando o limite
  \begin{align*}
  	\lim_{n\to\infty} S_n & = \lim \dfrac{\frac{9}{100}\left[1-\left(\frac{1}{100}\right)^n\right]}{1-\frac{1}{100}} \\[2mm]
  	                      & = \dfrac{\frac{9}{100}}{\;\frac{100-1}{100}\;}
                            = \dfrac{\;\frac{9}{100}\;}{\frac{99}{100}} \\[2mm]
  	                      & = \frac{9}{100}\;\frac{100}{99}
                            = \frac{9}{99}
                            = \frac{1}{11}
  \end{align*}
\end{answer}
\end{exercise}

% Seção 10.2 - Exercício 1, 3-6 página 20 Thomas
%------------------------------------------------------------------------------%
\begin{exercise}
  Encontre a fórmula para a $n$-ésima soma parcial de cada série
  e use-a para encontrar a soma da série se ela convergir.
  \begin{mathlist}
  \item 2 + \frac{2}{3}
          + \frac{2}{9}
          + \frac{2}{27}
          + \cdots
          + \frac{2}{3^{n-1}}
          + \cdots
  \item 1 - \frac{1}{2}
          + \frac{1}{4}
          - \frac{1}{8}
          + \cdots
          + \frac{(-1)^{n-1}}{2^{n-1}}
          + \cdots
  \item 1 - 2 + 4 - 8 + \cdots + (-1)^{n-1} 2^{n-1} + \cdots
  \item \frac{1}{2\cdot 3}
      + \frac{1}{3\cdot 4}
      + \frac{1}{4\cdot 5}
      + \cdots \) \\[2mm]
      \( \hspace{15.5mm}
      + \frac{1}{(n+1)\cdot (n+2)}
      + \cdots \)
  \item
    \frac{5}{1 \cdot 2} +
    \frac{5}{2 \cdot 3} +
    \frac{5}{3 \cdot 4} +
    \cdots +
    \frac{5}{n \cdot (n+1)} +
    \cdots
  \end{mathlist}
\end{exercise}

% Exercício 13 página 20 - Thomas
%------------------------------------------------------------------------------%
\begin{exercise}
Escreva os primeiros termos da série para
mostrar como a série começa. Então, calcule a soma da série.
\[
  \sum_{n=0}^{\infty} \frac{(-1)^n}{5^n}
\]
\begin{answer}
  Denotando a série por $S$ temos
  \begin{align*}
  	S & = \sum_{n=0}^{\infty} \frac{(-1)^n}{5^n}   \\
  	  & = \frac{(-1)^0}{5^0} + \frac{(-1)^1}{5^1} + \frac{(-1)^2}{5^2} + \cdots \\
  	  & = 1                  - \frac{1}{5}        + \frac{1}{25} + \cdots
  \end{align*}
  Observamos que essa série é uma série geométrica, com $\abs{r}=\sfrac{1}{5} < 1$,
  podemos então usar a formula da soma, porém precisamos primeiro corrigir as potências e índices
  \[
  	S = \sum_{n=0}^{\infty} \frac{(-1)^n}{5^n}
  	  = \sum_{n=1}^{\infty}\left(\frac{-1}{5}\right)^{n-1}
  \]
  Temos $\alpha=1$ e $r=\sfrac{-1}{5}$ e portanto
  \[
    \sum_{n=1}^{\infty} \left(\frac{-1}{5}\right)^{n-1}
    = \frac{1}{1-(\sfrac{-1}{5})}
    = \frac{5}{6}
  \]
\end{answer}
\end{exercise}

% Exercício 7-12, 14 página 20 - Thomas
%------------------------------------------------------------------------------%
\begin{exercise}
  Escreva os primeiros termos de cada série para
  mostrar como a série começa. Então, calcule a soma da série.
  \begin{mathlist}[2]
  \item \sum_{n=0}^{\infty} \frac{(-1)^n}{4^n}
  \item \sum_{n=2}^{\infty} \frac{1}{4^n}
  \item \sum_{n=1}^{\infty} \frac{7}{4^n}
  \item \sum_{n=0}^{\infty} (-1)^n \frac{5}{4^n}
  \item \sum_{n=0}^{\infty} \left(\frac{2^{n+1}}{5^n}\right)
  \end{mathlist}
\end{exercise}

% Exercício 35 página 20 - Thomas
%--------------------------------------------------------------------------------------------------%
\begin{exercise}
Encontre uma fórmula para a $n$-ésima soma parcial da série
e use-a para determinar se a série converge ou diverge.
Se a série converge, encontre a soma.
\[
  \sum_{n=1}^{\infty} \left( \frac{1}{n} -\frac{1}{n+1} \right)
\]
\begin{answer}
Expandindo os temos das somas somas parciais percebemos que esse é uma série telescópica
\begin{align*}
	S_n & = \sum_{k=1}^{n} \left( \frac{1}{k} -\frac{1}{k+1} \right) \\
	    & = \left( \frac{1}{1} -\frac{1}{2}\right)
        + \left( \frac{1}{2} -\frac{1}{3}\right) + \cdots \\
      & \qquad
        + \left( \frac{1}{n} -\frac{1}{n+1}\right) \\
	    & = 1 - \frac{1}{n+1}
\end{align*}
Calculando o limite temos
\[
  \lim_{n\to\infty} S_n = \lim_{n\to\infty} \left( 1 - \frac{1}{n+1} \right) = 1
\]
portanto a série converge para 1.
\end{answer}
\end{exercise}

% Exercício 36-40 página 20 - Thomas
%--------------------------------------------------------------------------------------------------%
\begin{exercise}
  Encontre uma fórmula para a $n$-ésima soma parcial de cada série
  e use-a para determinar se a série converge ou diverge.
  Se a série converge, encontre sua soma.
  \begin{mathlist}
  \item \sum_{n=1}^{\infty} \left( \frac{3}{n^2} - \frac{3}{(n+1)^2}\right)
  \item \sum_{n=1}^{\infty} \left( \ln\sqrt{n+1} - \ln\sqrt{n} \right)
  \item \sum_{n=1}^{\infty} \big( \tan(n) - \tan(n-1) \big)
  \item \sum_{n=1}^{\infty} \left( \cos^{-1}\left( \frac{1}{n+1}\right)
                                 - \cos^{-1}\left( \frac{1}{n+2}\right) \right)
  \item \sum_{n=1}^{\infty} \left( \sqrt{n+4} - \sqrt{n+3}\right)
  \end{mathlist}
\end{exercise}


% Exercício de prova do prof. André Pereira
%------------------------------------------------------------------------------%
\begin{exercise}
  Calcule a soma da série telescópica
  \[
    \sum_{n=1}^\infty \frac{2n+1}{n^2(n+1)^2}
  \]
\end{exercise}

% Exercício de prova do prof. André Pereira
%------------------------------------------------------------------------------%
\begin{exercise}
  Considere a série geométrica
  \[
    \sum_{n=10}^\infty \frac{(-9)^{(n-9)}}{10^{(n-10)}}
  \]
  \vspace{-\baselineskip}
  \begin{alphalist}
    \item Rearranje os índices da série geométrica acima para que ela fique na forma
    \(\sum_{n=1}^\infty ar^{n-1}\)
    e identifique $\alpha$ e $r$.
    \item Encontre o limite da série caso ela convirja.
  \end{alphalist}
\end{exercise}

% 69-78 Thomas pg 21
%------------------------------------------------------------------------------%
\begin{exercise}
  Escreva os primeiros termos das séries geométricas e encontre $\alpha$ e $r$.
  Em seguida determine para quais valores de $x$ temos que $\abs{r}<1$
  e calcule a soma de série.
  \begin{mathlist}
    % 69-71
    \item \sum_{n=0}^\infty (-1)^n x^n
    \item \sum_{n=0}^\infty (-1)^n x^{2n}
    \item \sum_{n=0}^\infty 3\left(\frac{x-1}{2}\right)^n
    % 73-78
    \item \sum_{n=0}^\infty 2^n x^n
    \item \sum_{n=0}^\infty (-1)^n x^{-2n}
    \item \sum_{n=0}^\infty (-1)^n (x+1)^n
    \item \sum_{n=0}^\infty \left(-\frac{1}{2}\right)^2 (x-3)^n
    \item \sum_{n=0}^\infty \sin^n(x)
    \item \sum_{n=0}^\infty \big(\ln(x)\big)^n
  \end{mathlist}
\end{exercise}

%------------------------------------------------------------------------------%
%\begin{exercise}
%\begin{answer}
%\end{answer}
%\end{exercise}

%------------------------------------------------------------------------------%
%\begin{exercise}
%\begin{mathlist}[2]
%  \item
%  \item
%  \item
%\end{mathlist}
%\begin{answer}
%  \begin{alphalist}[3]
%    \item
%    \item
%    \item
%  \end{alphalist}
%\end{answer}
%\end{exercise}

\end{exercises}
%------------------------------------------------------------------------------%
